\section{JADE --- adaptacyjna ewolucja różnicowa} \label{sec:jade}

Pierwszym powszechnie przyjętym rozwiązaniem problemu doboru parametrów sterujących DE był algorytm JADE (\textit{Adaptive Differential Evolution with Optional External Archive}). Został on zaproponowany przez Zhanga i Sandersona \cite{Zhang2009JADE} i wprowadza strategię mutacji \texttt{DE/current-to-pbest/1}:
\[
\mathbf{v}_{i,G} = \mathbf{x}_{i,G} + F_i \cdot (\mathbf{x}_{\mathrm{pbest},G} - \mathbf{x}_{i,G}) + F_i \cdot (\mathbf{x}_{r_1,G} - \tilde{\mathbf{x}}_{r_2,G})
\]
gdzie $\mathbf{x}_{\mathrm{pbest},G}$ to losowo wybrany osobnik spośród $p \cdot 100\%$ najlepszych w populacji, a $\tilde{\mathbf{x}}_{r_2,G}$ pochodzi z rozszerzonej puli, łączącej populację z \textit{archiwum}, czyli zbiorem wektorów wyeliminowanych w poprzednich pokoleniach \cite{Zhang2009JADE}. Archiwum zwiększa zróżnicowanie kierunków przeszukiwań bez konieczności rozbudowywania populacji.

Parametry $F$ i $CR$ w JADE są adaptowane indywidualnie dla każdego osobnika i są generowane odpowiednio z rozkładu Cauchy'ego i normalnego, wokół bieżących wartości średnich $\mu_F$ i $\mu_{CR}$. Po każdym pokoleniu wartości te są aktualizowane na podstawie parametrów osobników, którym udało się wygenerować lepsze potomstwo \cite{Zhang2009JADE}. Mechanizm ten pozwala algorytmowi samodzielnie dobierać wartości odpowiednie dla danego problemu i fazy przeszukiwań. JADE stał się punktem odniesienia dla kolejnych adaptacyjnych wariantów DE, a jego koncepcja archiwum i adaptacji parametrów jest omawiana w wielu pracach przeglądowych \cite{Bilal2020}.
